% !TEX root = ../termpaper.tex
% @author Leonhard Lüdeke
%

\section{Beitrag, Risiken und Arbeitsplan}

\subsection{Erwarteter wissenschaftlicher Beitrag}

Der erwartete wissenschaftliche Beitrag der Arbeit liegt in der systematischen Verbindung von Architekturkonzeption und Threat Modelling für ein LLM-basiertes Multi-Agenten-System im Unternehmenskontext.

Der Mehrwert besteht nicht in der Entwicklung eines neuen Agentenframeworks, sondern in einer nachvollziehbaren sicherheitsorientierten Entwurfsperspektive, die funktionale Komponenten, Vertrauensgrenzen, Risiken und Gegenmaßnahmen gemeinsam betrachtet.

Damit adressiert die Arbeit sowohl praktische als auch wissenschaftliche Aspekte.

Praktisch liefert sie eine strukturierte Grundlage für spätere Systementscheidungen im Unternehmenskontext.

Wissenschaftlich zeigt sie, wie sich klassische Konzepte der Multi-Agenten-Forschung und etablierte Sicherheitsüberlegungen auf neuere agentische Architekturen mit LLM-Komponenten übertragen und konkretisieren lassen.

\subsection{Annahmen, Unsicherheiten und offene Punkte}

Die Arbeit basiert auf mehreren Annahmen, die offen benannt werden müssen. Dazu zählen Annahmen über den vorgesehenen Nutzungskontext, über technische Integrationsmöglichkeiten, über interne und externe Kommunikationsbeziehungen sowie über potenzielle Angreifer und Missbrauchsszenarien.

Es ist wichtig, solche Annahmen transparent zu machen, da Bedrohungsbilder stark von Systemgrenzen, Rollenmodellen und Kommunikationsmustern abhängen.

Gleichzeitig bestehen Unsicherheiten, etwa hinsichtlich späterer Implementierungsdetails, konkreter Schnittstellen, Freigabeprozesse oder real verfügbarer Infrastruktur.

Diese Unsicherheiten mindern nicht den Wert der Arbeit, müssen jedoch explizit als Begrenzung des Geltungsbereichs berücksichtigt werden.